\documentclass{article}
\usepackage[utf8]{inputenc}
\usepackage{color}
\usepackage{url}

\newcommand{\red}[1]{\textcolor{red}{#1}}

\title{Lab Session 4 - Scientific Writing}
\author{Quang Hieu Trân}
\date{September 2024}

\begin{document}

\maketitle

\section{Introduction}
This lab follows the scientific writing course at Duke university. The exercises are also extracted from the course website.

You should revise the following sentences for clarity using the information explained in the following website:
\begin{itemize}
    \item \url{https://cgi.duke.edu/web/sciwriting/index.php}
\end{itemize}

\red{For each sentence you should justify your answer.}

\section{Subjects and Actions}
%Principles:
% 1. Put actions in verbs
% 2. Put characters in subjects
% 3. Keep subjects near verbs

\begin{enumerate}
        
    \item We subjected yeast to 20 min of amino acid deprivation and made ribosome footprint and mRNA-abundance measurements.
    \item[] \red{We deprived yeast of amino acids for 20 minutes and measured ribosome footprints and mRNA abundance.}
    \item[] \red{Instead of using the word "subjected," which can feel a bit formal and clunky, we just said what we did: "deprived yeast of amino acids" and "measured." It makes the sentence flow better and keeps things simple.}
    
    \item The number of different mechanisms that may exist for cells to interpret morphogens, and the importance of design features such as feedback or local cell-cell communication, is unclear.
    \item[] \red{It is unclear how many mechanisms exist for cells to interpret morphogens, and how important design features like feedback or local cell communication are.}
    \item[] \red{The original sentence was overloaded with information. By breaking it down and using "it is unclear," the message is much easier to follow. It also makes the sentence more straightforward, so you can focus on what matters without getting bogged down in complicated wording.}
     
    \item Localized fluctuations in substitution rate are widely employed to draw inference concerning the phenotypic significance of genomic sequence.
    \item[] \red{Researchers widely use localized fluctuations in substitution rates to infer the phenotypic significance of genomic sequences.}
    \item[] \red{Adding "Researchers widely use" makes it clear who’s doing the action. Instead of sounding vague and passive, the sentence is now more active and direct, making it easier to understand who’s making these inferences about the data.}
     
    \item In this study, we subjected seven different primate species to comparative analysis of the radial distribution pattern of human chromosome 18- and 19-homologous chromatin by three-dimensional fluorescence in situ hybridization.
    \item[] \red{In this study, we analyzed the radial distribution pattern of chromatin homologous to human chromosomes 18 and 19 in seven primate species using three-dimensional fluorescence in situ hybridization.}
    \item[] \red{We swapped out the wordy phrase "subjected ... to comparative analysis" with the simpler "analyzed," making the sentence cleaner and easier to read. It keeps the focus on what was done without making it sound too complicated.}
\end{enumerate}

\section{Cohesion, Coherence, and Emphasis}
%Principles:
% 1. Put new information last
% 2. Use passive voice judiciously
% 3. Make sure the first and last sentences of a paragraph match

\begin{enumerate}
    \item These data were then subjected to hierarchical clustering using the Pearson correlation as the distance metric.
    \item[] \red{We applied hierarchical clustering to these data, using Pearson correlation as the distance metric.}
    \item[] \red{Instead of the passive voice, we switched to active by starting with "We" to make it clear who’s doing the clustering. It makes the sentence flow better and helps the reader understand who’s responsible for the action while keeping the focus on the content.}
    
    \item We designated the putative GRB target genes as the human TFs that were under a HCNE density peak and whose orthologs were in conserved synteny with an array of HCNEs in human:zebrafish alignments. Density peaks and syntenic regions were downloaded from the Ancora Genome Browser.
    \item[] \red{We identified the putative GRB target genes as human TFs located under HCNE density peaks, whose orthologs were conserved with HCNE arrays in human alignments. Density peaks and syntenic regions were downloaded from the Ancora Genome Browser.}
    \item[] \red{We swapped "designated" with "identified" to make it clearer and more straightforward. Adding "We" at the start makes the sentence more active. By breaking the long sentence into two shorter ones, it’s easier to read and follow the main points.}
    
    \item Improvements are expected in the predictive power of all the scores being computed on multispecies alignments.
    \item[] \red{We expect improvements in the predictive power of all scores computed on multispecies alignments.}
    \item[] \red{Saying "We expect" makes it much clearer and puts the focus on the people doing the action. The sentence is also now active, which makes it more direct and easier to understand who's anticipating the improvements.}

    \item With the development of molecular embryology and the coming of the post-genomic era, the molecular mechanisms of morphological evolution have recently begun to be elucidated.
    \item[] \red{With the development of molecular embryology and the postgenomic era, researchers have recently begun elucidating the molecular mechanisms of morphological evolution.}
    \item[] \red{By using "researchers," we clarify who’s doing the work and eliminate the vagueness of the passive voice. It’s now more obvious who’s taking action and what’s being clarified, making the sentence more direct and to the point.}
    
\end{enumerate}

\section{Concision and Simplicity}
% Principles:
%  1. Omit needless words
%  2. Prefer simple words
%  3. Use simple subjects
%  4. Use adjectives frugally
\begin{enumerate}
    
    \item One of the complexes formed by the hematopoietic transcription factor Gata1 is a complex with the Ldb1 and Tal1 proteins
    \item[] \red{The hematopoietic transcription factor Gata1 forms a complex with Ldb1 and Tal1 proteins.}
    \item[] \red{By cutting out the repetitive part "one of the complexes... is a complex," the sentence gets straight to the point. It's much shorter and easier to follow now.}
    
    \item A phylogenetic tree based on the number of differentially expressed genes between species recapitulates their known phylogeny.
    \item[] \red{A phylogenetic tree, constructed from differentially expressed genes between species, accurately reflects their known phylogeny.}
    \item[] \red{In this rephrased sentence, we used "constructed from" to make the connection between the tree and the genes clearer. It ties everything together nicely, showing how the tree is built using the gene data. The sentence also flows better because it first tells you what the tree is and then what it does, which makes it easier to follow. Finally, by saying "accurately reflects," we’re putting emphasis on the accuracy of the tree in representing the species’ relationships.}
    
    \item It is important to note for some of our later discussion that we permit a block to have only one row. 
    \item[] \red{We allow a block to have only one row, which is important for later discussion.}
    \item[] \red{Cutting out "It is important to note" makes the sentence less cluttered. The sentence is now more focused and highlights the key information without any extra fluff.}
    
    \item These sequences are generally among the least constrained in the genome, so they constitute a plausible neutral standard accounting for regional variation in mutation and recombination rates.
    \item[] \red{These sequences are minimally constrained, making them a plausible neutral standard for regional variation in mutation and recombination rates.}
    \item[] \red{Instead of saying "are generally among the least constrained," we just went with "are minimally constrained." It’s much more direct, and cutting out the extra bits makes the whole sentence clearer and easier to understand.}
    
\end{enumerate}

\end{document}
